\cdcpartie{Contexte et définition du besoin}

\cdcsection{Contexte et définition du problème}

\cdcsubsection{Contexte du projet}

Au Maroc, la réparation de téléphones est assurée en grande partie par de petites boutiques de quartier. Leur activité repose encore largement sur le papier ou sur la mémoire du réparateur pour suivre les appareils reçus, les réparations en cours, les paiements et les coordonnées des clients.

Le projet Qribna consiste à développer une plateforme web qui digitalise la gestion de ces boutiques et qui facilite, pour les clients, la recherche d’un réparateur adapté à leur problème. Il est réalisé par une équipe de cinq étudiants dans le cadre du module de génie logiciel.

\cdcsubsection{Définition du problème}

Le problème se pose à la fois du côté des boutiques et du côté des clients.

\cdcintertitre{Du côté des boutiques}

De nombreuses petites boutiques gèrent encore les appareils reçus, les réparations, les paiements et les informations des clients sur papier ou de mémoire. Cette organisation entraîne plusieurs difficultés.

\begin{cdcitems}
  \item Des pertes d’informations et des retards dans le traitement des réparations.
  \item Des conflits avec les clients sur l’état initial des appareils, faute de trace à leur réception.
  \item Un manque de visibilité sur les revenus de la boutique, les paiements et les commissions.
\end{cdcitems}

\cdcintertitre{Du côté des clients}

Les clients, de leur côté, ne disposent d’aucun outil pour choisir et suivre une réparation.

\begin{cdcitems}
  \item Il leur est difficile de trouver une boutique adaptée à leur problème.
  \item Ils ne connaissent pas le prix estimé d’une réparation avant de se déplacer.
  \item Ils ne peuvent pas suivre l’avancement des travaux sans appeler ou passer à la boutique.
  \item Ils ne peuvent pas toujours choisir un moyen de paiement pratique.
\end{cdcitems}

\cdcsubsection{Solution proposée}

Qribna est une plateforme unique qui relie les clients, les réparateurs et les propriétaires de boutiques. Elle réunit trois dimensions.

\begin{cdcitems}
  \item \textbf{La gestion professionnelle des boutiques} avec l’enregistrement des appareils, la photo et la validation de l’état initial, le suivi des statuts, les paiements et le calcul des commissions.
  \item \textbf{Une expérience simplifiée pour les clients} avec la recherche géolocalisée selon la panne, le diagnostic préliminaire, la comparaison des boutiques, le suivi par WhatsApp et le choix du moyen de paiement.
  \item \textbf{L’analyse du marché marocain} grâce à un observatoire des prix et des tendances de la réparation, fondé sur des données agrégées.
\end{cdcitems}

Les éléments qui distinguent Qribna sont le suivi des réparations par WhatsApp, la validation photographique de l’état initial, la recherche géolocalisée selon le type de panne, le diagnostic préliminaire avec estimation du prix, le paiement en espèces ou électronique, la gestion des commissions, le suivi des garanties, l’observatoire statistique des prix et une interface adaptée aux habitudes locales.
